\section{THỰC TRẠNG HOẠT ĐỘNG QUẢN LÝ TẠI FPT}

\subsection{Sơ đồ cơ cấu tổ chức bộ máy quản lý}
Mô hình tổ chức của FPT kết hợp giữa cơ cấu chức năng tập trung tại Tập đoàn mẹ và cơ cấu ma trận theo đơn vị kinh doanh chiến lược (SBU) tại các công ty con.
% Chèn sơ đồ cơ cấu tổ chức vào đây
% \begin{figure}[htbp]
%     \centering
%     \includegraphics[width=0.85\textwidth]{figures/so_do_co_cau_quan_ly.png}
%     \caption{Sơ đồ bộ máy điều hành và quản lý Tập đoàn FPT}
%     \label{fig:co_cau_quan_ly}
% \end{figure}

\subsection{Bộ máy điều hành: Phân quyền Tổng Giám đốc và các Phó Tổng Giám đốc}
\begin{itemize}
    \item \textbf{Tổng Giám đốc (CEO - Nguyễn Văn Khoa):} Chịu trách nhiệm trước HĐQT về toàn bộ kết quả hoạt động sản xuất kinh doanh thường nhật, trực tiếp điều hành việc hiện thực hóa các chiến lược 3 năm và kế hoạch năm.
    \item \textbf{Các Phó Tổng Giám đốc chức năng:}
    \begin{itemize}
        \item \textit{Phó TGĐ phụ trách Chuyển đổi số:} Điều phối năng lực công nghệ, chuẩn hóa hệ thống dữ liệu toàn tập đoàn.
        \item \textit{Phó TGĐ phụ trách Tài chính / Giám đốc Tài chính (CFO):} Quản trị dòng tiền, phân bổ vốn và thẩm định dự án đầu tư.
        \item \textit{Giám đốc Công nghệ (CTO):} Định hướng R\&D công nghệ mới, sản phẩm Made by FPT.
    \end{itemize}
\end{itemize}

\subsection{Cơ chế vận hành và phối hợp mẹ - con}
FPT sử dụng hệ thống quản trị mục tiêu thông qua OKR (Objectives and Key Results) và BSC/KPI để đồng bộ hóa mục tiêu từ cấp tập đoàn xuống từng đơn vị thành viên và nhân viên.

\subsection{Điểm nghẽn và hạn chế cần cải thiện trong công tác quản lý}
\begin{enumerate}
    \item \textbf{Bộ máy phình to, nhiều tầng nấc trung gian:} Với quy mô vượt 70.000 nhân sự trên toàn cầu, độ trễ truyền đạt mệnh lệnh và thủ tục phê duyệt hành chính nội bộ gia tăng, giảm tính linh hoạt trước biến động thị trường.
    \item \textbf{Xung đột quyền lực giữa thế hệ sáng lập và lãnh đạo điều hành trẻ:} Nhóm lãnh đạo cốt cán sáng lập vẫn giữ ảnh hưởng lớn đến các quyết định tác nghiệp cụ thể, tạo rào cản tâm lý cho đội ngũ điều hành thế hệ mới.
    \item \textbf{Silo dữ liệu giữa các công ty con:} Các đơn vị thành viên (FPT Telecom, FPT Software, FPT IS) hoạt động tương đối độc lập, dẫn đến hiện tượng cát cứ thông tin, chưa tận dụng tối đa giá trị dữ liệu chéo của khách hàng.
\end{enumerate}

\subsection{Đề xuất giải pháp hoàn thiện cơ chế quản lý}
\begin{itemize}
    \item Tái cấu trúc theo hướng phẳng hóa bộ máy quản lý trung gian.
    \item Phân cấp triệt để quyền tự chủ tài chính và nhân sự cho các đơn vị kinh doanh chiến lược.
    \item Xây dựng nền tảng dữ liệu dùng chung (Data Lakehouse) toàn tập đoàn để phá vỡ các silo dữ liệu.
\end{itemize}
